\section{Evaluation}
\label{sec:evaluation}


\subsection{Evaluation Settings}

\begin{figure*}[t!]
    \centering
    \begin{subfigure}[t]{0.32\textwidth}
        \centering
        \includegraphics[width=\linewidth]{Figures/evaluation/overall_lvbench.pdf}
        \caption{LVBench}
        \label{subfig:overall_performance_lvbench}
    \end{subfigure}
    \hfill
    \begin{subfigure}[t]{0.32\textwidth}
        \centering
        \includegraphics[width=\linewidth]{Figures/evaluation/overall_videomme.pdf}
        \caption{VideoMME-Long}
        \label{subfig:overall_performance_videomme}
    \end{subfigure}
    \hfill
    \begin{subfigure}[t]{0.32\textwidth}
        \centering
        \includegraphics[width=\linewidth]{Figures/evaluation/overall_ava100.pdf}
        \caption{\sysname{}-100}
        \label{subfig:overall_performance_ava100}
    \end{subfigure}
    \vspace{-1mm}
    \caption{The achieved accuracy of \sysname{} and various baselines on the LVBench, VideoMME-Long, and \sysname{}-100 benchmarks.}
    \label{fig:overall_acc}
    \vspace{-2mm}
\end{figure*}

\subsubsection{Benchmarks}
\sysname{} is evaluated on two widely used public long-video benchmarks and one ultra-long video benchmarks proposed by us, covering a broad range of video scenarios and problem types.

\textbf{LVBench}\cite{wang2024lvbench} stands out among publicly available benchmarks for its exceptionally long average video duration, approximately 4100 seconds per video. It comprises 103 videos with a total of 1549 questions, covering six distinct video domains and addressing six task types including temporal grounding, summarization, and reasoning.

\textbf{VideoMME-Long}\cite{fu2024videomme} is a subset of the VideoMME benchmark, focusing on videos exceeding 20 minutes in duration, with an average length of 2400 seconds. Comprising a total of 300 videos and 900 questions, the benchmark covers a wide range of video themes, spanning 6 primary visual domains with 30 subfields to ensure broad scenario generalizability, and includes 12 distinct task types.

\textbf{\sysname{}-100} is proposed by us, which is an ultra-long video benchmark specially designed to evaluate video analysis capabilities
\sysname{}-100 consists of 8 videos, each exceeding 10 hours in length, and includes a total of 120 manually annotated questions. The benchmark
covers four typical video analytics scenarios: human daily activities, city walking, wildlife monitoring, and traffic monitoring, each scenario contains two videos. The human daily activity scenario features egocentric videos selected and stitched from the Ego4D\cite{grauman2022ego4d}. City walking and wildlife monitoring videos are curated from publicly available recordings on YouTube, capturing urban exploration and animal monitoring respectively. Traffic monitor videos are composed from clips in the Bellevue Traffic Video Dataset~\cite{bhardwaj2022ekya}. All questions are carefully designed by human annotators, who also provide reference answers as the ground truth. In addition, GPT-4o is utilized to generate plausible distractor options. The accuracy is evaluated by analyzing \sysname's responses to multiple-choice questions included in the benchmarks. \revised{More details regarding \sysname-100 are provided in \S\ref{app:benchmark}.}



\subsection{Baselines}
We conduct a comprehensive comparison between \sysname{} and a wide range of baseline models, encompassing both mainstream VLMs and specialized Video-RAG methods. The VLM baselines include GPT-4o~\cite{achiam2023gpt}, Gemini-1.5-Pro~\cite{team2023gemini}, Phi-4-Multimodal~\cite{abdin2024phi4}, Qwen2.5-VL-7B~\cite{bai2025qwen2_5}, InternVL2.5-8B~\cite{chen2024internvl}, and LLaVA-Video-7B~\cite{zhang2024llavavideo}. Each of these models is evaluated with two typical strategies: uniform sampling and vectorized retrieval, where a CLIP-based retriever selects the top-K relevant frames based on the user query. In addition to VLMs, we benchmark \sysname{} against SOTA Video-RAG frameworks, including VideoTree\cite{wang2024videotree}, VideoAgent\cite{wang2024videoagent}, DrVideo\cite{ma2024drvideo}, and VCA\cite{yang2024vca}. Among these, VideoTree, VideoAgent, and VCA are built upon GPT-4o, while DrVideo leverages GPT-4.

\begin{figure}[!h]
    \centering
    \includegraphics[width=0.85\linewidth]{Figures/evaluation/task_specific_performance.pdf}
    \caption{The accuracy achieved by \sysname{} and the baselines across typical query categories on LVBench: Temporal Grounding (TG), Summarization (SU), Reasoning (RE), Entity Recognition (ER), Event Understanding (EU), and Key Information Retrieval (KIR).\vspace{-0.25cm}}
    \label{fig:task_specific}
\end{figure}

\subsection{Overall Evaluation}

\subsubsection{Overall Performance}

Fig.~\ref{fig:overall_acc} illustrates the overall accuracy achieved by \sysname{} compared to various baselines on the LVBench, VideoMME-Long, and \sysname{}-100 benchmarks. Across all three benchmarks, \sysname{} consistently outperforms the baselines. Specifically, on LVBench, \sysname{} delivers a remarkable 16.9\% improvement, while on VideoMME-Long, it advances the SOTA by approximately 5.2\%. On the \sysname{}-100 benchmark, \sysname{} achieves an accuracy of 75.8\%, significantly surpassing all competing methods.

In detail, compared to video-RAG methods, \sysname{} achieves improvements of 21\% and 7.8\% on LVBench and VideoMME-Long, respectively. When compared to vectorized retrieval-based methods, \sysname{} demonstrates gains of 16.9\% on LVBench and 20.8\% on \sysname{}-100. Furthermore, against uniform sampling baselines, \sysname{} improves performance by approximately 19.6\% and 26.9\% on LVBench and \sysname{}-100, respectively.

Notably, on \sysname{}-100, when evaluated with extremely long videos, \sysname{} maintains robust performance, whereas the baselines degrade significantly. This highlights the effectiveness of \sysname{} in handling L4 video analytics tasks.

\subsubsection{Performance on Different Query Categories}
We also evaluate the accuracy achieved by \sysname{} across typical query categories on LVBench. As illustrated in Fig.~\ref{fig:task_specific}, our approach achieves improvements of 16\%, 5.3\%, 35.6\%, 21.2\%, 17.5\%, and 18.9\% across six key task types: Temporal Grounding, Summarization, Reasoning, Entity Recognition, Event Understanding, and Key Information Retrieval, respectively, compared to the uniform sampling and vectorized retrieval baselines powered by Gemini-1.5-Pro. Notably, \sysname{} demonstrates particularly strong performance on reasoning tasks, which require identifying causal relationships between events and linking preceding and succeeding events within the video. This highlights \sysname{}'s ability to effectively locate and extract critical information from long videos, thereby enabling advanced L4 video analytics systems.
\begin{figure}[!t]
    \centering
    \includegraphics[width=0.85\linewidth]{Figures/evaluation/overall_combined.pdf}
    \caption{The accuracy achieved by \sysname{} and baselines across three benchmarks when utilizing different LLMs and VLMs.}
    \label{fig:ava_configs}
\end{figure}

\subsubsection{Performance under Different Configurations}
Fig.~\ref{fig:ava_configs} shows the performance of \sysname{} using different models configurations for SA and CA. For SA, two models were used: Qwen2.5 14B and 32B. For CA, two models were used: Qwen2.5-VL-7B and Gemini-1.5-Pro. 
The results show that across the three benchmarks, \sysname{} using Gemini-1.5-Pro for CA achieved improvements of 18.9\%, 5.2\%, and 20.8\% respectively compared to the best baseline result using the same model, while using Qwen2.5-VL-7B yielded improvements of 13\%, 7.2\%, and 15\% respectively, fully demonstrating the effectiveness of our method. Notably, even when only using Qwen2.5-32B and Qwen2.5-7B based on the textual content from EKG without accessing raw frames, \sysname{} can surpass the performance of Qwen2.5-VL-7B on the three benchmarks and also outperform most models shown in Figs.~\ref{subfig:overall_performance_lvbench}, \ref{subfig:overall_performance_videomme}, and \ref{subfig:overall_performance_ava100}.

\subsubsection{Performance on Different Video Lengths}

\begin{figure}[t]
    \centering
    \resizebox{0.8\linewidth}{!}{\includegraphics{Figures/evaluation/videomme_concat.pdf}}
    \caption{The accuracy achieved by \sysname{} and the baselines across varying video lengths via concatenating videos from LVBench.}
    \label{fig:concat}
\end{figure}



To evaluate the robustness of \sysname{} with respect to video length, we conducted experiments on videos of varying durations. Specifically, sequences of 3.3, 6.6, and 10 hours were created by concatenating videos from the VideoMME-Long benchmark. Performance was measured using identical questions across these varying video lengths. As illustrated in Fig.~\ref{fig:concat}, both Qwen2.5-VL-7B and Gemini-1.5-Pro baselines exhibit significant performance degradation as video length increases. When extended to 10 hours, their performance declines by 4.6\% and 8.2\%, respectively, under the uniform sampling method, compared to the original VideoMME-Long benchmark. For the vectorized retrieval setting, the performance drops are 4.6\% and 5.5\%, respectively. These results highlight the limitations of these methods in scaling effectively with increasing video length. In contrast, \sysname{} consistently maintains stable performance across all video lengths, underscoring its robustness and scalability in handling video data of any duration.


\subsubsection{System Overhead}




\sysname{} is designed to enable the near real-time construction of EKGs. As shown in Fig.~\ref{fig:construction overhead}, we measured the average processing speed (in FPS) of \sysname{} while constructing EKGs from LVBench videos across various hardware platforms, with the input video stream fixed at 2 FPS. On 2 $\times$ A100 GPUs, \sysname{} achieved an impressive processing speed of 6.7 FPS, significantly exceeding the input stream rate. On a single RTX 4090, a typical edge server hardware, \sysname{} maintained a processing speed of 4.4 FPS, still well above the input frame rate. Even on a single RTX 3090, \sysname{} performed effectively, achieving 2.5 FPS. This performance demonstrates its capability to support efficient, near real-time EKG construction for L4 video analytics. The overhead during the retrieval and generation phases will be discussed further in \S\ref{subsec:ablation}.





\begin{figure}[t]
    \centering
    \resizebox{0.85\linewidth}{!}{\includegraphics{Figures/evaluation/construction_overhead.pdf}}
    \caption{Total index construction overhead evaluated on various types of typical edge server hardware.}
    \label{fig:construction overhead}
\end{figure}


\revised{
To evaluate system overhead during the generation phase, we conducted a detailed latency analysis of three stages within the generation pipeline, as summarized in Table~\ref{tab:generation_overhead}. The experiment is performed on one A100 GPU with LMDeploy~\cite{2023lmdeploy}.~\footnote{\revised{We configure $cache\_max\_entry\_count$ to 0.3 in LMDeploy, restricting the KV cache memory consumption to no more than 30\% of the total GPU memory.}} The retrieval stage utilizing JinaCLIP demonstrates notable efficiency, requiring only 0.44 seconds per query and consuming less than 1 GB of GPU memory on average. The agentic search stage presents the predominant source of latency: 101.5 seconds with Qwen2.5-14B, rising to 174.2 seconds with the larger Qwen2.5-32B, thereby underscoring the inference time of LLMs as the principal runtime overhead. During the consistency-enhanced generation stage, Qwen2.5-VL-7B incurs a latency of 45.8 seconds and approximately 31 GB of memory usage, whereas the API-based Gemini-1.5-Pro achieves substantially lower latency at 14.2 seconds. The results demonstrate that the agentic searching stage represents the principal performance bottleneck, underscoring the importance of strategic model selection to balance accuracy and computational efficiency. Furthermore, the results also highlight the imperative of optimizing the agentic searching process, which we identify as a key direction for future research.
}

\begin{table}[!t]
\centering
\resizebox{0.9\linewidth}{!}{
\begin{tabular}{cccc}
\toprule[1.5pt]
\textbf{Stage} & \textbf{Model} & \textbf{Latency (s)} & \textbf{GPU Memory (GB)} \\
\midrule
Tri-View Retrieval & JinaCLIP & 0.44 & 0.8 \\
\midrule
\multirow{2}{*}{Agentic Searching} & Qwen2.5-14B & 101.5 & 30 \\
 & Qwen2.5-32B & 174.2 & 40 \\
\midrule
\multirow{2}{*}{Consistency Enhanced Gen.} & Qwen2.5-VL-7B & 45.8 & 31 \\
 & Gemini-1.5-Pro & 14.2 & - \\ 
\bottomrule[1.5pt]
\end{tabular}
}
\caption{\revised{The breakdown of latency and GPU memory consumption across the three stages of the generation phase in \sysname. All measurements are performed on a single A100 GPU, with LLMs deployed via LMDeploy~\cite{2023lmdeploy} utilizing AWQ~\cite{lin2024awq}.}}
\label{tab:generation_overhead}
\end{table}

\subsection{Ablation Evaluation}
\label{subsec:ablation}
We randomly sampled 20 videos and 305 corresponding questions from LVBench for our ablation study. All ablation experiments were conducted on 2 $\times$ A100 GPUs.

\begin{table}[!t]
\centering
\label{tab:context_length}
\begin{tabular}{lcc} 
\toprule[1.5pt]
\textbf{Method} & \textbf{Acc.} & \textbf{Overhead(h)}  \\ 
\midrule
MiniRAG         &     28.1          &          3.49                   \\
LightRAG        &   30.6            &   3.52                          \\
\sysname{}         &      \textbf{39.7}         &       \textbf{0.31}                    \\
\bottomrule[1.5pt]
\end{tabular}
\caption{The achieved accuracy and construction overhead evaluated when using EKG and KG as index in \sysname{} and baseline models on the subset of LVBench. The total video duration is around 1.2 hours.}
\label{tab:ekg_kg}
\end{table}

\subsubsection{Different Index Construction Methods}
We compare \sysname{}’s EKG construction method with two representative knowledge graph-based 
construction methods: LightRAG~\cite{guo2024lightrag} and MiniRAG~\cite{fan2025minirag}. Since both of them only support text-only construction, we use the full set of descriptions obtained through the semantic chunking (\S\ref{subsec: sematic_chunking}) as their input textual corpus. 
We use Qwen2.5 7B to construct EKG and KG for \sysname and baselines, respectively. In the retrieval and generation phase, we use the same LLM, Qwen2.5 14B and the same settings, \eg maximum number of tokens of retrieved events or entities.


As shown in Table~\ref{tab:ekg_kg}, \sysname{} demonstrates a significant performance advantage over the baselines. Specifically, it achieves 11.6\% higher accuracy than MiniRAG and 9.1\% higher accuracy than LightRAG. Crucially, this improved performance comes with substantially less construction overhead, requiring only 0.31 hours compared to 3.49 and 3.52 hours for the baselines. The rationale is that baselines construct KG based on massive uniform chunks, while \sysname utilizes the semantic chunks. This substantial gap in both effectiveness and efficiency highlights that \sysname{}'s EKG construction method not only yields higher-quality knowledge representations but also drastically reduces the time needed to build the graph.


\subsubsection{Different Tree Search Depths}


We also evaluate the achieved accuracy and overhead applying different settings in the retrieval phase, \ie tree search depth. The effect of tree depth lies in a crucial trade-off: while shallower depths may struggle to retrieve comprehensive information, increasing the depth allows access to richer information from deeper nodes. However, this comes with a significant increase in tree search overhead, and the information from deeper levels can introduce more noise, potentially negatively impacting the final generation quality. Table~\ref{tab:tree depth} presents the results comparing different tree search depths on performance and tree search overhead. As shown, performance generally increases with increasing tree depth up to a certain point. Specifically, for all three \sysname{} configurations evaluated, the highest accuracy is achieved at a tree search depth of 3. Accuracy decreases when the depth is further increased to 4, suggesting that excessive depth leads to the retrieval of detrimental noise or irrelevant information, outweighing the benefit of additional context. Conversely, the tree search overhead increases sharply with depth. Expanding the search from depth 1 (6.7s) to depth 2 (27.3s) incurs a moderate increase. 
Comparing the accuracy improvements and the overhead increase, a tree search depth of 3 offers the optimal balance. 

\label{subsub:depth}
\begin{table}[t!]
\centering
\resizebox{\linewidth}{!}{
\begin{tabular}{lcccc} 
\toprule[1.5pt]
\multicolumn{1}{c}{\multirow{2}{*}{\textbf{Method}}} & \multicolumn{4}{c}{\textbf{Tree Search Depth}}  \\
\multicolumn{1}{c}{}                                 & 1    & 2    & 3             & 4                 \\
\midrule
\sysname{}(Qwen2.5 14B)                                   & 34.1 & 36.1 & \textbf{40.9} & 39.5              \\
\sysname{}(Qwen2.5 14B + Qwen2.5VL 7B)                    & 49.3 & 52.1 & \textbf{53.8} & 50.2              \\
\sysname{}(Qwen2.5 14B + Gemini-1.5-Pro)                  & 54.2 & 58.4 & \textbf{61.5} & 52.7              \\
\midrule[1pt]
\multicolumn{1}{c}{\textbf{Tree Search Overhead(s)}}               & 6.7 & 27.3 & 90.1 & 370.3 \\
\bottomrule[1.5pt]
\end{tabular}
}
\caption{The achieved accuracy and overhead when applying different tree search depths in the agentic search of \sysname evaluated on the subset of LVBench.}
\vspace{-1em}
\label{tab:tree depth}
\end{table}






\subsubsection{Different Consistency Evaluation Settings}
For consistency-enhanced generation, \sysname{} incorporates two key parameters: $\lambda$, which governs the balance between the contributions of thought consistency and answer consistency, and the number of generations for self-consistency evaluation. Fig.~\ref{subfig:self_consistency_weight} illustrates the impact of varying $\lambda$ values on the accuracy achieved by \sysname{}. Notably, the optimal performance is observed when $\lambda$ is set to $0.3$, highlighting the importance of jointly considering both intermediate thought consistency and final answer consistency to ensure robust results. As depicted in Fig.~\ref{subfig:self_consistency_times}, the accuracy of \sysname{} gradually improves as the number of self-consistency iterations increases. However, this improvement comes at the expense of significantly higher computational overhead. For example, increasing the self-consistency iterations from 8 to 16 yields only a 0.9\% accuracy gain, while nearly doubling the computational cost. This demonstrates a clear trade-off between marginal accuracy improvements and resource efficiency. Balancing this trade-off, we adopt 8 self-consistency iterations in the implementation of \sysname{}, ensuring a practical balance between performance and computational overhead.

